\HMTNarrativeHeading{Persistencia y trayectoria electrónica}

La estructura electrónica empieza en la posición central persistente
de APP. La inversión celular distingue esa posición de las parejas
conjugadas que la rodean. Su evaluación conserva las hojas, los
residuos y los cocientes; la trayectoria prolonga esos datos mediante
las operaciones del TPK. La exposición de la masa recibe después las
secciones de acción y los lectores que utiliza.

La vacancia central permite ver una diferencia decisiva entre lectura
residual e historia aritmética. Dos posiciones pueden compartir la
suma y el residuo multiplicativo, mientras sus productos enteros
conservan cocientes distintos. El cambio de cociente registra la
transferencia; la orientación distingue sus posiciones conjugadas.
El ejemplo hace visible qué información necesita una reconstrucción
que restituya la operación realizada.

La cadena electrónica añade una cronología explícita. La recurrencia
celular, las inversiones y las ventanas de lectura conservan los
registros de la sección persistente. El calendario registra los pasos
en que permanece la capacidad; el defecto central registra la diferencia
de cociente entre evaluaciones con igual residuo. La trayectoria
electrónica conserva este segundo dato junto con su orientación y su
cronología. La vacancia de capacidad y el defecto de cociente cuentan
operaciones distintas, con sus respectivos registros.

Esta continuidad vuelve pertinente la distinción entre retorno y
recuperación. Una fase puede repetirse mientras el registro de la
sección incorpora el trayecto. La evaluación posterior utiliza esa
información transportada y conserva identificada la sección central.
La masa, la orientación y las expresiones de acción se presentan con
sus antecedentes, de modo que cada paso pueda leerse como una
composición determinada.

El electrón ofrece así un caso desarrollado de la tesis del artículo.
La persistencia corresponde a una estructura; la trayectoria describe
su evolución; los registros distinguen los cambios; el lector
dimensional obtiene una magnitud. Su interés para la conservación
reside en esa sucesión completa. La unidad que se mantiene durante
el transporte conserva relaciones cuyo contenido se hace accesible
en lecturas diferentes. El estudio operatorio posterior abstraerá
las condiciones de recuperación y mostrará sus balances exactos.
